\section{Kiểm thử và kết quả}

\subsection{Phạm vi kiểm thử}
Do thời hạn 6 tuần chỉ dành trọn một tuần cho kiểm thử tích hợp
và hệ thống, nhóm tập trung kiểm thử vào các luồng nghiệp vụ có rủi ro
cao nhất, thay vì phủ toàn bộ 38 use case:
\begin{itemize}
    \item Xác thực \& phân quyền (RBAC).
    \item Đăng ký hoạt động có kiểm soát sức chứa (race-condition).
    \item Check-in bằng QR/RFID/NFC.
    \item Ghi nhận giao dịch ngoại tuyến và đồng bộ khi có mạng trở lại.
    \item Đối soát giao dịch và xuất hóa đơn cuối chuyến.
\end{itemize}

\textbf{Không nằm trong phạm vi kiểm thử của phiên bản 1.0:} Module BLE
định vị hành khách (chưa hiện thực), kiểm thử tải với quy mô trên 1.000
hành khách đồng thời (do giới hạn hạ tầng thử nghiệm).

\subsection{Chiến lược kiểm thử}
\begin{itemize}
    \item \textbf{Unit Test:} Jest, mục tiêu tỷ lệ pass tối thiểu 90\%
    cho các service xử lý business rule quan trọng (kiểm tra sức chứa,
    tính số dư tài khoản).
    \item \textbf{Integration Test:} Supertest kiểm thử API kết hợp
    PostgreSQL/Redis thật trên môi trường staging.
    \item \textbf{System Test:} Kiểm thử thủ công theo kịch bản
    (test case) trên Web Admin, Mobile App và giả lập thiết bị POS mất
    mạng.
\end{itemize}

\subsection{Kết quả kiểm thử}

\begin{table}[htbp]
    \centering
    \small
    \caption{Tổng hợp kết quả kiểm thử}
    \begin{tabularx}{\textwidth}{|Y|p{2cm}|p{2cm}|p{2cm}|Y|}
        \hline
        \textbf{Module} & \textbf{Số test case} & \textbf{Pass} & \textbf{Fail} & \textbf{Ghi chú} \\
        \hline
        Xác thực \& RBAC & -- & -- & -- & Điền số liệu thực tế sau khi
        chạy test \\
        \hline
        Đăng ký hoạt động & -- & -- & -- & \\
        \hline
        Check-in QR/RFID/NFC & -- & -- & -- & \\
        \hline
        Giao dịch ngoại tuyến \& đồng bộ & -- & -- & -- & Module rủi ro
        cao nhất, cần kiểm thử kỹ nhất \\
        \hline
        Đối soát \& hóa đơn & -- & -- & -- & \\
        \hline
    \end{tabularx}
\end{table}

\textit{Ghi chú: Bảng trên để trống cột số liệu — nhóm điền kết quả thật
sau khi chạy xong bộ kiểm thử, đính kèm chi tiết đầy đủ tại file
Report5\_Test Report.xls theo quy ước bàn giao ở Chương 2.}

\subsection{Lỗi phát hiện và hướng xử lý}
Các lỗi phát hiện trong quá trình test được ghi nhận theo mức độ ưu tiên
(Critical/High/Medium/Low) trên công cụ theo dõi lỗi (GitHub Issues),
tuân theo quy trình Quản lý chất lượng đã mô tả ở Chương 2 — lỗi Critical
liên quan tới đồng bộ giao dịch (sai lệch số dư tài khoản) được ưu tiên
xử lý trong vòng 24 giờ.